% Options for packages loaded elsewhere
\PassOptionsToPackage{unicode}{hyperref}
\PassOptionsToPackage{hyphens}{url}
\PassOptionsToPackage{dvipsnames,svgnames,x11names}{xcolor}
%
\documentclass[
  11pt,
  a4paper,
]{article}
\usepackage{amsmath,amssymb}
\usepackage{iftex}
\ifPDFTeX
  \usepackage[T1]{fontenc}
  \usepackage[utf8]{inputenc}
  \usepackage{textcomp} % provide euro and other symbols
\else % if luatex or xetex
  \usepackage{unicode-math} % this also loads fontspec
  \defaultfontfeatures{Scale=MatchLowercase}
  \defaultfontfeatures[\rmfamily]{Ligatures=TeX,Scale=1}
\fi
\usepackage{lmodern}
\ifPDFTeX\else
  % xetex/luatex font selection
\fi
% Use upquote if available, for straight quotes in verbatim environments
\IfFileExists{upquote.sty}{\usepackage{upquote}}{}
\IfFileExists{microtype.sty}{% use microtype if available
  \usepackage[]{microtype}
  \UseMicrotypeSet[protrusion]{basicmath} % disable protrusion for tt fonts
}{}
\makeatletter
\@ifundefined{KOMAClassName}{% if non-KOMA class
  \IfFileExists{parskip.sty}{%
    \usepackage{parskip}
  }{% else
    \setlength{\parindent}{0pt}
    \setlength{\parskip}{6pt plus 2pt minus 1pt}}
}{% if KOMA class
  \KOMAoptions{parskip=half}}
\makeatother
\usepackage{xcolor}
\usepackage[margin=24mm]{geometry}
\usepackage{longtable,booktabs,array}
\usepackage{calc} % for calculating minipage widths
% Correct order of tables after \paragraph or \subparagraph
\usepackage{etoolbox}
\makeatletter
\patchcmd\longtable{\par}{\if@noskipsec\mbox{}\fi\par}{}{}
\makeatother
% Allow footnotes in longtable head/foot
\IfFileExists{footnotehyper.sty}{\usepackage{footnotehyper}}{\usepackage{footnote}}
\makesavenoteenv{longtable}
\setlength{\emergencystretch}{3em} % prevent overfull lines
\providecommand{\tightlist}{%
  \setlength{\itemsep}{0pt}\setlength{\parskip}{0pt}}
\setcounter{secnumdepth}{-\maxdimen} % remove section numbering
\usepackage{amsmath,amssymb}
\usepackage{microtype}
\usepackage{xurl}
\usepackage{fancyhdr}
\usepackage{longtable,booktabs,array}
\setlength{\emergencystretch}{2em}
\setlength{\parindent}{0pt}
\setlength{\parskip}{0.45em}
\pagestyle{fancy}
\fancyhf{}
\fancyhead[L]{\small Alternating Motzkin sums}
\fancyhead[R]{\small Research note | 28 October 2026}
\fancyfoot[C]{\thepage}
\setlength{\headheight}{14pt}
\renewcommand{\headrulewidth}{0.3pt}
\allowdisplaybreaks[2]
\ifLuaTeX
  \usepackage{selnolig}  % disable illegal ligatures
\fi
\usepackage{bookmark}
\IfFileExists{xurl.sty}{\usepackage{xurl}}{} % add URL line breaks if available
\urlstyle{same}
\hypersetup{
  pdftitle={Sharp log-concavity for alternating Motzkin sums},
  pdfauthor={Dev D. Goyal},
  colorlinks=true,
  linkcolor={Black},
  filecolor={Maroon},
  citecolor={Blue},
  urlcolor={Black},
  pdfcreator={LaTeX via pandoc}}

\title{Sharp log-concavity for alternating Motzkin sums}
\usepackage{etoolbox}
\makeatletter
\providecommand{\subtitle}[1]{% add subtitle to \maketitle
  \apptocmd{\@title}{\par {\large #1 \par}}{}{}
}
\makeatother
\subtitle{A bounded-forcing theorem, exact parameter thresholds, and an
auditable proof package}
\author{Dev D. Goyal}
\date{28 October 2026}

\begin{document}
\maketitle

\par\noindent{\normalsize\bfseries This work is fully autonomous agentic research. GPT 6.1 Sol in Codex
produced formal proofs and verification software for the accompanying
package. Dev D. Goyal is responsible for publication.\par}

\textbf{Status.} The bounded-forcing Theorem 2 and its corrected
scaled-forcing corollary have complete Lean proofs. The
concrete-sequence, uniform-family, root-transfer and
threshold-classification claims have conventional proofs and are outside
that formalization. A public manuscript dated 17 October 2026 already
claims the original root-log-concavity conjecture. No first-priority
claim is made for its resolution or for the extensions below. A broader
novelty assessment remains open.

\section{1. Definitions and results}\label{definitions-and-results}

Define the Motzkin numbers by \(M_0=M_1=1\) and \[
(n+3)M_{n+1}=(2n+3)M_n+3nM_{n-1}\qquad(n\ge1),
\] and put \[
T_n=\sum_{k=0}^n(-1)^{n-k}M_k.
\] Zhao {[}Z, Conjecture 1, p.~8{]} asks whether \((T_n^{1/n})_{n\ge9}\)
is log-concave. A prior public manuscript {[}P{]} claims strict
log-concavity on that tail. The argument here instead uses a rational
invariant interval and exact polynomial identities. The general passage
from ratio log-concavity to root log-concavity is established prior
methodology {[}CGW{]}; the elementary multiplicative identity is
included to make the argument self-contained, without claiming that
principle as new.

A positive sequence \((b_n)_{n\ge N}\) is \emph{strictly log-concave}
when \(b_n^2>b_{n-1}b_{n+1}\) at every center \(n\ge N+1\). In
particular, the statement that \((a_{n+1}/a_n)_{n\ge N}\) is strictly
log-concave means \[
a_na_{n+2}^3>a_{n+1}^3a_{n+3}\qquad(n\ge N).
\] These indexing conventions are important for sharpness.

\textbf{Theorem 1 (original sequence).} For all integers \(n\ge10\), \[
T_nT_{n+2}^{3}>T_{n+1}^{3}T_{n+3}.
\tag{1}
\] For all integers \(n\ge6\), \[
T_n^{n^2-1}>
T_{n-1}^{n(n+1)/2}T_{n+1}^{n(n-1)/2}.
\tag{2}
\] Consequently the quotient tail starts at 10 and the root tail starts
at 5. Both indices are optimal. In addition, for every \(n\ge12\), \[
T_nT_{n+2}^{3}-T_{n+1}^{3}T_{n+3}
>\frac{800T_n^4}{(n+3)^3}.
\tag{3}
\] The constant 800 is a convenient certified constant, not an asserted
optimum.

The next result covers arbitrary bounded recurrence errors. Define \[
L_n=3-\frac9{2(n+3)}+\frac9{16(n+3)^2},\qquad
U_n=L_n+\frac2{(n+3)^3},\qquad E_n=\frac1{8(n+3)^2}.
\tag{4}
\]

\textbf{Theorem 2 (bounded forcing).} Let real numbers \(a_n\),
\(n\ge12\), satisfy \[
a_{12}>1800,\qquad L_{12}<\frac{a_{13}}{a_{12}}<U_{12},
\tag{5}
\] and, for every integer \(m\ge14\), \[
(m+2)a_m=m(2a_{m-1}+3a_{m-2})+\varepsilon_m,
\qquad |\varepsilon_m|\le2.
\tag{6}
\] Then for all \(n\ge12\), \[
a_n>8(n+3)^2,\qquad L_n<\frac{a_{n+1}}{a_n}<U_n,
\tag{7}
\] and \[
a_na_{n+2}^3-a_{n+1}^3a_{n+3}
>\frac{800a_n^4}{(n+3)^3}>0.
\tag{8}
\] There is no alternation or independence hypothesis on the errors.

\textbf{Corollary (scaled forcing).} Let \(B\ge0\). Replace the error
bound in (6) by \(|\varepsilon_m|\le B\), and assume \(a_{12}>900B\)
with the same strict quotient seed in (5). Then, for every \(n\ge12\),
\[
a_n>4B(n+3)^2,\qquad L_n<\frac{a_{n+1}}{a_n}<U_n,
\] and the quartic gap in (8) holds with the same constant 800. At
\(B=0\), the first conclusion is \(a_n>0\).

For \(B>0\), apply Theorem 2 to \(2a_n/B\). Its growth bound becomes
\(a_n>4B(n+3)^2\), and division by the fourth power of the scale
preserves (8). For \(B=0\), multiply the sequence by one sufficiently
large positive constant so that its initial value exceeds 1800. Theorem
2 gives positivity and the quotient interval; cancellation of the same
scale gives (8).

For the family extension, define the Riordan numbers by \(R_0=1\),
\(R_1=0\), and \[
(n+1)R_n=(n-1)(2R_{n-1}+3R_{n-2})\qquad(n\ge2).
\] This is the standard sequence {[}O1{]}. Put \[
A_n(s)=R_{n+1}+s(-1)^n=T_n+(s-1)(-1)^n,
\qquad 0\le s\le1.
\tag{9}
\] The equality in (9) is proved below. The endpoints are shifted
Riordan numbers and the original alternating Motzkin sums.

\textbf{Theorem 3 (uniform family).} Every \(A_n(s)\) with \(n\ge2\) is
positive. Uniformly for \(0\le s\le1\), the quotient inequality holds
for \(n\ge10\), the root inequality holds at centers \(n\ge6\), and the
quantitative inequality (8) holds for \(n\ge12\), with \(a_n=A_n(s)\).
The best tail indices for each individual parameter are completely
classified in Section 6.

\section{2. The inhomogeneous
recurrence}\label{the-inhomogeneous-recurrence}

From the definition, \[
M_n=T_n+T_{n-1}\qquad(n\ge1).
\] For \(n\ge2\), let \[
D_n=(n+2)T_n-n(2T_{n-1}+3T_{n-2}).
\] Substitution in the Motzkin recurrence at index \(n-1\) gives, for
\(n\ge3\), \[
\begin{aligned}
D_n+D_{n-1}
={}&(n+2)(T_n+T_{n-1})\\
&-(2n+1)(T_{n-1}+T_{n-2})
-3(n-1)(T_{n-2}+T_{n-3})=0.
\end{aligned}
\] The initial values give \(D_2=2\). Therefore \[
(n+2)T_n=n(2T_{n-1}+3T_{n-2})+2(-1)^n
\qquad(n\ge2).
\tag{10}
\] Adding \((s-1)(-1)^n\) changes the forcing term to \(2s(-1)^n\). At
\(s=0\) the resulting homogeneous recurrence, and its first two values 0
and 1, identify it with \(R_{n+1}\) by induction. This proves (9) and \[
(n+2)A_n(s)=n(2A_{n-1}(s)+3A_{n-2}(s))+2s(-1)^n.
\tag{11}
\]

The finite values used throughout are \[
\begin{array}{c|rrrrrrrr}
n&0&1&2&3&4&5&6&7\\\hline
T_n&1&0&2&2&7&14&37&90
\end{array}
\] \[
\begin{array}{c|rrrrrrrr}
n&8&9&10&11&12&13&14&15\\\hline
T_n&233&602&1586&4212&11299&30536&83098&227474.
\end{array}
\] They follow directly from the defining recurrence. The independent
audit also regenerates them from \[
M_n=\sum_{j=0}^{\lfloor n/2\rfloor}\binom n{2j}\frac1{j+1}\binom{2j}j.
\] This second construction is a computational cross-check, not an
additional hypothesis of the proof.

\section{3. Proof of the bounded-forcing
theorem}\label{proof-of-the-bounded-forcing-theorem}

\subsection{3.1. Invariant interval and
growth}\label{invariant-interval-and-growth}

For \(n\ge12\) one has \[
L_n-\frac83=\frac{16n^2-120n-477}{48(n+3)^2}>0,
\qquad
3-U_n=\frac{72n^2+423n+589}{16(n+3)^3}>0.
\tag{12}
\] For the first numerator, writing \(n=12+t\) gives \(16t^2+264t+387\).
Thus \(8/3<L_n<U_n<3\).

Let \[
f_n(r,e)=\frac{n+1}{n+3}\left(2+\frac3r\right)+\frac{2e}{n+3}.
\] For positive \(r\), this is decreasing in \(r\) and increasing in
\(e\). These facts follow by subtraction of reciprocal fractions and
require no differentiation.

The two assertions in (7) follow by joint induction. They hold at
\(n=12\) by (5). Suppose they hold at \(k-1\), where \(k\ge13\). Since
\(q_{k-1}=a_k/a_{k-1}>8/3>2\), \[
a_k>16(k+2)^2>8(k+3)^2;
\] the last inequality is \(2(k+2)^2-(k+3)^2=k^2+2k-1>0\). The
recurrence at \(m=k+1\) gives \[
q_k=f_k(q_{k-1},e_k),\qquad e_k=\frac{\varepsilon_{k+1}}{2a_k},
\qquad |e_k|<E_k.
\] Exact rational identities give \[
f_n(U_{n-1},-E_n)-L_n
=\frac{304n^3+1687n^2+1599n-1790}
{16(n+3)^3(48n^3+216n^2+297n+146)},
\tag{13}
\] \[
U_n-f_n(L_{n-1},E_n)
=\frac{112n^2-509n-1331}
{16(n+3)^3(16n^2+40n+19)}.
\tag{14}
\] All denominators are positive. At \(n=13+t\), the two numerators
become \[
304t^3+13543t^2+199589t+971988,
\qquad112t^2+2403t+10980,
\] so both are positive for \(n\ge13\). Monotonicity of \(f_n\) now
proves \(L_k<q_k<U_k\), completing the joint induction.

\subsection{3.2. A local algebraic
comparison}\label{a-local-algebraic-comparison}

Fix \(n\ge12\). For independent real variables \(r,u,v\), put \[
y=\frac{(n+2)(2r+3)+2u}{n+4},\quad
z=\frac{(n+3)(2y+3r)+2v}{n+5},\quad
F_n(r,u,v)=y^3-r^3z.
\tag{15}
\] At \[
r=q_n,\qquad u=\frac{\varepsilon_{n+2}}{2a_n},\qquad
v=\frac{\varepsilon_{n+3}}{2a_n},
\] the recurrence implies \(y=a_{n+2}/a_n\) and \(z=a_{n+3}/a_n\).
Consequently \[
\frac{a_na_{n+2}^3-a_{n+1}^3a_{n+3}}{a_n^4}=F_n(r,u,v).
\tag{16}
\] The invariant interval gives \(8/3<r<3\) and
\(|u|,|v|<E_n\le1/1800<1/100\).

On the larger rectangular region \[
8/3\le r\le3,\qquad |u|,|v|\le1/100,
\tag{17}
\] one has \(7\le y\le9\) and \(z\ge19\). Indeed, the lower bound on
\(y-7\) is \[
\frac{200n-1703}{150(n+4)}>0,
\] the lower bound on \(9-y\) is \(899/[50(n+4)]>0\), and, using
\(y\ge7\), the lower bound on \(z-19\) is \[
\frac{150n-1451}{50(n+5)}>0.
\] Define the slopes \[
\alpha=\frac{2(n+2)}{n+4},\qquad
\beta=\frac{(n+3)(7n+20)}{(n+4)(n+5)}.
\] Then \(0<\alpha\le2\) and \(\beta\ge5\), the latter because \[
(n+3)(7n+20)-5(n+4)(n+5)=2n^2-4n-40>0.
\]

For \(r\le t\) with \(u,v\) fixed, write \(y_r,z_r\) and \(y_t,z_t\) for
the corresponding values. Since \(y_t-y_r=\alpha(t-r)\) and
\(z_t-z_r=\beta(t-r)\), \[
\begin{aligned}
F_n(t,u,v)-F_n(r,u,v)
=(t-r)\bigl[&\alpha(y_t^2+y_ty_r+y_r^2)\\
&-(t^2+tr+r^2)z_r-t^3\beta\bigr].
\end{aligned}
\tag{18}
\] The bracket is at most \[
486-3(8/3)^2\,19-(8/3)^3\,5=-\frac{382}{27}<0.
\] Thus \(F_n\) decreases in its first variable. The checker
additionally verifies explicit nonnegative decompositions underlying
this bound.

For \(u_1\le u_2\) with \(r,v\) fixed, put \[
\gamma=\frac2{n+4},\qquad \eta=\frac{2(n+3)}{n+5}\le2.
\] Direct factorization gives \[
F_n(r,u_2,v)-F_n(r,u_1,v)
=(u_2-u_1)\gamma\bigl[y_2^2+y_2y_1+y_1^2-\eta r^3\bigr].
\tag{19}
\] The bracket is at least \(147-54=93>0\), so \(F_n\) increases in its
second variable. Finally, \[
F_n(r,u,v_2)-F_n(r,u,v_1)=-\frac{2r^3(v_2-v_1)}{n+5},
\tag{20}
\] so it decreases in its third variable. All intermediate comparison
points remain in (17). Hence \[
F_n(q_n,u,v)\ge F_n(U_n,-E_n,E_n).
\tag{21}
\]

\subsection{3.3. The corner certificate and quantitative
gap}\label{the-corner-certificate-and-quantitative-gap}

At the right-hand corner of (21), exact substitution gives \[
F_n(U_n,-E_n,E_n)
=\frac{P(n)}{65536(n+3)^{11}(n+4)^3(n+5)}.
\tag{22}
\] All coefficients of \(P\) are positive; they are listed in Appendix
A. Here is a compact independently checkable form of the identity. Set
\(m=n+3\) and \[
\begin{aligned}
D&=16m^3,& V&=48m^3-72m^2+9m+32,\\
Y&=(n+2)(2V+3D)-4m,\\
Z&=m\bigl(2Y+3V(n+4)\bigr)+4m(n+4).
\end{aligned}
\] At the corner, \(r=V/D\), \(y=Y/[D(n+4)]\), and
\(z=Z/[D(n+4)(n+5)]\). The polynomial identity \[
(n+5)DY^3-(n+4)^2V^3Z=(n+3)P(n)
\tag{23}
\] implies (22), since \(D^4=65536(n+3)^{12}\).

For the explicit constant in (8), the second certificate is \[
P(n)-800\cdot65536(n+3)^8(n+4)^3(n+5)
=\sum_{j=0}^{12}c_j(n-12)^j,
\tag{24}
\] where every \(c_j\) is strictly positive. Appendix A gives the exact
coefficients. Equation (24) implies that (22) is strictly larger than
\(800/(n+3)^3\) for every \(n\ge12\). Combining (16), (21), and (22)
proves Theorem 2. No limit, numerical optimization, or finite-range
extrapolation is used.

\section{4. Application to the parameter
family}\label{application-to-the-parameter-family}

The two seeds are \[
A_{12}(s)=11298+s,\qquad A_{13}(s)=30537-s.
\] For \(0\le s\le t\le1\), \[
\frac{30537-s}{11298+s}-\frac{30537-t}{11298+t}
=\frac{41835(t-s)}{(11298+s)(11298+t)}\ge0.
\] The endpoint comparisons are exact: \[
\frac{1081}{400}<\frac{30536}{11299},\qquad
\frac{30537}{11298}<\frac{145967}{54000},
\tag{25}
\] where \(L_{12}=1081/400\) and \(U_{12}=145967/54000\). The respective
cross-product gaps are 181 and 137166. Also \(11298>1800\). Equations
(11) and (25) therefore satisfy Theorem 2 for every parameter in
\([0,1]\). This proves the quotient and quantitative claims for all
\(n\ge12\).

The values \(A_n(s)\) for \(2\le n\le13\) are positive directly from the
finite table in Section 2. Theorem 2 supplies positivity at all larger
indices.

Let \[
Q_n(s)=A_n(s)A_{n+2}(s)^3-A_{n+1}(s)^3A_{n+3}(s).
\tag{26}
\] For \(2\le n\le11\), the quartic terms cancel and the resulting cubic
coefficients are listed in Appendix B. The even-index cubics are
strictly increasing on \([0,1]\), and the odd-index cubics are strictly
decreasing. A purely algebraic check is as follows. If
\(p(s)=c_0+c_1s+c_2s^2+c_3s^3\), then \[
\frac{p(t)-p(s)}{t-s}=c_1+c_2(s+t)+c_3(s^2+st+t^2)
\quad(s<t).
\] For each even index, \(c_3>0\), \(c_2<0\), and \(c_1+2c_2>0\). For
each odd index, \(c_3<0\) and \(c_1+2\max(c_2,0)<0\). These facts prove
the asserted strict monotonicity.

In particular, \[
Q_{10}(s)\ge Q_{10}(0)=2279044915731>0,
\] \[
Q_{11}(s)\ge Q_{11}(1)=59137804243970>0.
\] This completes the uniform quotient statement for every \(n\ge10\).

\section{5. Algebraic transfer to root
log-concavity}\label{algebraic-transfer-to-root-log-concavity}

For a positive sequence at the relevant indices, define \[
K_n=\frac{a_n^{n^2-1}}
{a_{n-1}^{n(n+1)/2}a_{n+1}^{n(n-1)/2}}
\qquad(n\ge2).
\] Both half-exponents are integers. The root-log-concavity inequality
at center \(n\) is equivalent to \(K_n>1\): raise the root inequality to
the positive power \(n(n^2-1)/2\).

Cancellation of monomials gives the exact identity \[
\frac{K_{n+1}}{K_n}
=\left(\frac{a_{n-1}a_{n+1}^3}{a_n^3a_{n+2}}\right)^{n(n+1)/2}.
\tag{27}
\] For example, the exponents of \((a_{n-1},a_n,a_{n+1},a_{n+2})\) on
both sides agree after expanding the four quadratic polynomials in
\(n\). The checker verifies these exponent identities in the
doubled-exponent form, which avoids division by 2.

For \(a_n=A_n(s)\), the quotient statement proved above makes the
right-hand side of (27) larger than one whenever \(n\ge11\). It remains
to check \(K_n>1\) for \(6\le n\le11\). For a fixed center \(n\), the
root gap \[
H_n(s)=A_n(s)^{n^2-1}
-A_{n-1}(s)^{n(n+1)/2}A_{n+1}(s)^{n(n-1)/2}
\tag{28}
\] is strictly increasing in \(s\) when \(n\) is even, and strictly
decreasing when \(n\) is odd. Indeed the center and its two neighbors
vary in opposite directions, all bases are positive, and all exponents
are positive. Hence the worst endpoints are 0 for even centers and 1 for
odd centers. The six exact checks are \[
\begin{array}{rcl}
36^{35}&>&15^{21}91^{15},\\
90^{48}&>&37^{28}233^{21},\\
232^{63}&>&91^{36}603^{28},\\
602^{80}&>&233^{45}1586^{36},\\
1585^{99}&>&603^{55}4213^{45},\\
4212^{120}&>&1586^{66}11299^{55}.
\end{array}
\tag{29}
\] These are integer comparisons, checked without floating point.
Equation (27) propagates the last one to every later center. This proves
the uniform root statement, and hence Theorem 3 apart from the sharp
classification.

At \(s=1\), the immediate preceding obstructions are \[
Q_9(1)=-92071225688<0,
\qquad14^{24}<7^{15}37^{10}.
\tag{30}
\] The latter follows, after canceling \(7^{15}\), from \[
2^{24}7^9=677021181018112
<4808584372417849=37^{10}.
\] Therefore the original quotient tail cannot start before 10 and the
original root tail cannot start before 5. This proves all claims in
Theorem 1.

A further quantitative consequence of (8) and \(q_j<3\) is available for
any positive sequence satisfying Theorem 2. Whenever \(n\ge13\), \[
\frac{K_{n+1}}{K_n}>
\left(1+\frac{800}{729(n+2)^3}\right)^{n(n+1)/2}.
\tag{31}
\] Indeed, the fraction inside (27) is \(1+Q_{n-1}/(a_n^3a_{n+2})\), and
\(a_n^3a_{n+2}<729a_{n-1}^4\). No claim is made that (31) alone supplies
a positive base root gap for arbitrary initial seeds.

\section{6. Exact sharp thresholds for every
parameter}\label{exact-sharp-thresholds-for-every-parameter}

\subsection{6.1. Quotient tails}\label{quotient-tails}

For \(j\in\{4,5,6,7,9\}\), let \(\rho_j\) be the unique zero of
\(Q_j(s)\) in \((0,1)\). Strict monotonicity from Section 4 and the
rational sign brackets in Appendix B establish existence and uniqueness.
The order is \[
\rho_6<\rho_4<\rho_5<\rho_7<\rho_9.
\tag{32}
\] Define \(N_q(s)\) to be the least positive integer such that the
quotient tail \((A_{n+1}(s)/A_n(s))_{n\ge N_q(s)}\) is defined,
positive, and strictly log-concave. Then \[
\begin{array}{c|c}
\text{Parameter interval}&N_q(s)\\\hline
{[0,\rho_6]}&7\\
(\rho_6,\rho_4]&5\\
(\rho_4,\rho_5)&3\\
{[\rho_5,\rho_7)}&6\\
{[\rho_7,\rho_9)}&8\\
{[\rho_9,1]}&10
\end{array}
\tag{33}
\] Here and below, equality at a boundary counts as failure of
\emph{strict} log-concavity.

To prove (33), note first that \(Q_8(s)>0\) throughout \([0,1]\), since
\(Q_8\) is increasing and \(Q_8(0)=68892649>0\). Also all \(Q_n\) with
\(n\ge10\) are positive. The signs of \(Q_4,Q_5,Q_6,Q_7,Q_9\) now follow
from their unique roots and monotonicity. The last nonpositive indices
in the six rows of (33) are respectively \(6,4,2,5,7,9\).

Only the third row requires an additional low-index observation. The
isolating interval for \(\rho_5\) lies below \(6/25\). Direct
substitution gives \(Q_2(6/25)<0\) and \(Q_3(6/25)>0\). Since \(Q_2\) is
increasing and \(Q_3\) is decreasing, \(Q_2(s)<0<Q_3(s)\) throughout
\((\rho_4,\rho_5)\). Thus the listed last nonpositive index really is 2.
This proves every case and every endpoint in (33).

\subsection{6.2. Root tails}\label{root-tails}

For \(j=2,3,4,5\), let \(\tau_j\) be the unique zero in \((0,1)\) of
\(H_j(s)\) from (28). The four defining polynomials may be given without
expansion: \[
\begin{aligned}
H_2(s)&=(1+s)^3-(1-s)^3(3-s),\\
H_3(s)&=(3-s)^8-(1+s)^6(6+s)^3,\\
H_4(s)&=(6+s)^{15}-(3-s)^{10}(15-s)^6,\\
H_5(s)&=(15-s)^{24}-(6+s)^{15}(36+s)^{10}.
\end{aligned}
\tag{34}
\] The parity monotonicity argument applies also at these indices; at
the endpoint \(s=1\), zero values such as \(A_1(1)=0\) do not disturb
the polynomial sign argument. The brackets in Appendix B prove that each
polynomial changes sign, and \[
\tau_4<\tau_2<\tau_3<\tau_5.
\tag{35}
\] There is a simple exact form for one boundary: \[
H_2(s)=-(s^2-6s+1)(s^2-s+2),\qquad \tau_2=3-2\sqrt2.
\tag{36}
\]

Let \(N_r(s)\) be the least positive integer such that the root tail
\((A_n(s)^{1/n})_{n\ge N_r(s)}\) is positive and strictly log-concave.
Then \[
\begin{array}{c|c}
\text{Parameter interval}&N_r(s)\\\hline
{[0,\tau_4]}&4\\
(\tau_4,\tau_2]&2\\
(\tau_2,\tau_3)&1\\
{[\tau_3,\tau_5)}&3\\
{[\tau_5,1]}&5
\end{array}
\tag{37}
\] All centers at least 6 are already strictly positive. In the five
rows, the last nonpositive centers among \(2,3,4,5\) are respectively
\(4,2\), none, \(3,5\). Their signs follow directly from (34)--(35) and
parity monotonicity. A last failing center \(j\) forces a root tail to
start at least at \(j\), not \(j+1\), because its first internal center
is one later than its starting index. In the middle row \(s<1\), so
\(A_1(s)=1-s>0\), and all four small-center inequalities hold. Thus
starting index 1 is valid and minimal. This proves (37).

The endpoint consequences are particularly simple: shifted Riordan
numbers have the sharp quotient-tail start 7 and root-tail start 4; the
alternating Motzkin sums have starts 10 and 5. The classification
concerns \(R_{n+1}^{1/n}\) at \(s=0\), not the differently indexed
sequence \(R_n^{1/n}\).

\section{7. Verification scope and remaining
work}\label{verification-scope-and-remaining-work}

The standard-library checker uses exact integer and rational polynomial
arithmetic. It compares coefficients and shifted positivity
certificates, regenerates finite data, and checks interval signs and
base cases. The recorded \texttt{python\ -O\ -S} run did not depend on
assertions or third-party packages.

A separate SymPy audit uses independently written expressions and finite
Motzkin values from Catalan and binomial formulas. Deliberate
certificate mutations test rejection. Finite parameter and index
regressions are diagnostics and establish no infinite claim.

An exact checker is not a theorem prover. The written induction and
comparisons link the identities to the universal conclusions.

The separate \texttt{formalization/} package proves Theorem 2 and the
corrected scaled-forcing corollary for arbitrary real sequences and
errors. It includes the simultaneous growth and quotient bounds and the
quartic margin with constant 800 for every \(n\ge12\). The scaled growth
is \(a_n>4B(n+3)^2\) for every \(B\ge0\), including positivity at
\(B=0\). The stated seeds and recurrence onset are unchanged.

The package README gives the pinned Lean and Mathlib builds and the
statement and axiom audits. Only \texttt{propext},
\texttt{Classical.choice}, and \texttt{Quot.sound} are permitted. Theorems 1 and 3,
concrete-sequence identification, root transfer and both threshold
classifications retain their conventional proofs and are outside this
formalization.

The original conjecture has a prior public proof claim. The priority and
significance of the extensions and parameter classifications remain
unresolved. No certified first solution is claimed.


\section{Appendix A. Integer polynomial
certificates}\label{appendix-a.-integer-polynomial-certificates}

Write \(P(n)=\sum_{j=0}^{12}p_jn^j\) and use \(c_j\) as in (24). Every
entry below is positive. The machine-readable source is
\texttt{proof/certificates.json}.

\begin{longtable}[]{@{}rrr@{}}
\toprule\noalign{}
\(j\) & \(p_j\) & \(c_j\) \\
\midrule\noalign{}
\endhead
\bottomrule\noalign{}
\endlastfoot
0 & 4660409556032 & 306994825721596657664 \\
1 & 24132003073520 & 706171963822972745984 \\
2 & 56628081108524 & 429272400422485948736 \\
3 & 79753800966621 & 133590768426096292509 \\
4 & 75163688028504 & 25822213571611442448 \\
5 & 49975719381810 & 3377331981471101346 \\
6 & 24049942325892 & 311625831704217912 \\
7 & 8443039471929 & 20633351214169401 \\
8 & 2146478817312 & 978753038000288 \\
9 & 385445807712 & 32566142496352 \\
10 & 46409331456 & 723458556672 \\
11 & 3364180992 & 9653801984 \\
12 & 111034368 & 58605568 \\
\end{longtable}

\section{Appendix B. Family cubics and exact root
brackets}\label{appendix-b.-family-cubics-and-exact-root-brackets}

For \(Q_n(s)=c_0+c_1s+c_2s^2+c_3s^3\), the coefficients are:

\begin{longtable}[]{@{}rrrrr@{}}
\toprule\noalign{}
\(n\) & \(c_0\) & \(c_1\) & \(c_2\) & \(c_3\) \\
\midrule\noalign{}
\endhead
\bottomrule\noalign{}
\endlastfoot
2 & -189 & 756 & -36 & 43 \\
3 & 2349 & -9504 & 54 & -102 \\
4 & -27189 & 134784 & -234 & 250 \\
5 & 479373 & -2074888 & -6 & -628 \\
6 & -4865265 & 34034060 & -2934 & 1608 \\
7 & 160155377 & -586941772 & -9186 & -4181 \\
8 & 68892649 & 10545295603 & -72309 & 11009 \\
9 & 103947123741 & -196017930893 & -389241 & -29295 \\
10 & 2279044915731 & 3749897573668 & -2548848 & 78655 \\
11 & 132662633568165 & -73524813169500 & -15941880 & -212815 \\
\end{longtable}

The following intervals are exact rational isolating intervals, with
common denominator \(10^6\). They are not floating-point root
approximations used as proof. At both endpoints the checker evaluates
the defining polynomial exactly, with the sign orientation determined by
monotonicity.

\begin{longtable}[]{@{}lrr@{}}
\toprule\noalign{}
Boundary & Lower numerator & Upper numerator \\
\midrule\noalign{}
\endhead
\bottomrule\noalign{}
\endlastfoot
\(\rho_6\) & 142954 & 142955 \\
\(\rho_4\) & 201778 & 201779 \\
\(\rho_5\) & 231031 & 231032 \\
\(\rho_7\) & 272862 & 272863 \\
\(\rho_9\) & 530293 & 530294 \\
\(\tau_4\) & 57244 & 57245 \\
\(\tau_2\) & 171572 & 171573 \\
\(\tau_3\) & 407203 & 407204 \\
\(\tau_5\) & 530571 & 530572 \\
\end{longtable}

\section{References and provenance}\label{references-and-provenance}

{[}Z{]} Feng-Zhen Zhao, \emph{The Log-Balancedness of the Sequence for
the Alternating Sums of Motzkin Numbers}, INTEGERS \textbf{26} (2026),
A55. Published 1 May 2026. DOI: 10.5281/zenodo.19949645.
\href{https://math.colgate.edu/~integers/aa55/aa55.pdf}{Journal-hosted
paper}. Conjecture 1 on page 8 is the original target; the paper's other
results are not claimed here as new.

{[}CGW{]} William Y. C. Chen, Jeremy J. F. Guo, and Larry X. W. Wang,
\emph{Infinitely log-monotonic combinatorial sequences}, Advances in
Applied Mathematics \textbf{52} (2014), 99--120. DOI:
10.1016/j.aam.2013.08.003.
\href{https://www.sciencedirect.com/science/article/pii/S0196885813000973}{Publisher
page}. Ratio-to-root transfer is prior methodology.

{[}P{]} Public manuscript, \emph{Motzkin alternating sums: strict
log-concavity of the root sequence} (Chinese title; English description
here), in the \texttt{xzwwwwww/math} repository. The relevant source was
added on 17 October 2026 in commit
\texttt{2e173134a5da24512ff9d2fbf1504b5634d29af8}.
\href{https://github.com/xzwwwwww/math/blob/2e173134a5da24512ff9d2fbf1504b5634d29af8/Motzkin\%E7\%8C\%9C\%E6\%83\%B3/motzkin_root_log_concavity.tex}{Pinned
source}. It claims strict root log-concavity at centers \(n\ge10\) by an
integral method. This package records that public claim and does not
independently certify every step of that manuscript or its publication
status.

{[}O1{]} OEIS Foundation, \href{https://oeis.org/A005043}{A005043:
Riordan numbers}. Used to identify the standard homogeneous sequence;
the necessary recurrence connection is proved in Section 2.

{[}O2{]} OEIS Foundation, \href{https://oeis.org/A187306}{A187306:
alternating sum of Motzkin numbers}. Sequence identification, not a
source of the infinite inequalities.

All literature status statements refer to the sources inspected on 28
October 2026. A negative search result is not proof of novelty.
Third-party full-text papers and manuscripts are not redistributed in
this package.

\end{document}
